\documentclass[a4paper]{article}
\usepackage{titlesec}
\usepackage{geometry}
\geometry{hmargin={2cm,2cm}, vmargin={3.5cm,3.5cm}}
\usepackage{graphicx}
\usepackage[table]{xcolor}
\usepackage{hyperref}
\hypersetup{pdfborder = {0 0 0},
colorlinks,
citecolor=black,
filecolor=black,
linkcolor=black,
urlcolor=blue}
\definecolor{logo}{RGB}{20,99,163}
\usepackage[T1]{fontenc}
\usepackage[utf8]{inputenc}
\usepackage[english,italian]{babel}
\usepackage{fancyhdr}
\usepackage{lastpage}
\usepackage{comment}

% creazione comando subsubsubsection
\titleclass{\subsubsubsection}{straight}[\subsection]

\newcounter{subsubsubsection}[subsubsection]
\renewcommand\thesubsubsubsection{\thesubsubsection.\arabic{subsubsubsection}}
\titleformat{\subsubsubsection}
    {\normalfont\normalsize\bfseries}{\thesubsubsubsection}{1em}{}
    \titlespacing*{\subsubsubsection}
    {0pt}{3.25ex plus 1ex minus .2ex}{1.5ex plus .2ex}

\makeatletter
\def\toclevel@subsubsubsection{4}
\def\l@subsubsubsection{\@dottedtocline{4}{7em}{4em}}
\makeatother

\setcounter{secnumdepth}{4}
\setcounter{tocdepth}{4}
%fine creazione comando


%Comandi di impaginazione uguale per tutti i documenti
\pagestyle{fancy}
\lhead{\includegraphics[scale=0.2]{../../../../Images/logo.png}}
%Titolo del documento
\rhead{\large Verbale Esterno del 2021-04-19}
%\rfoot{\thepage}
\cfoot{Pagina \thepage{} di \pageref{LastPage}}
\setlength{\headheight}{45pt}
\renewcommand{\footrulewidth}{0.4pt}
\begin{document}

\begin{titlepage}
    \begin{center}
        \includegraphics{../../../../Images/logo.png}\\
        \vspace{20px}
        \textcolor{logo}{\hrulefill}\\
        \vspace{20px}
        \textbf{\huge\textcolor{logo}{Verbale Esterno del 2021-04-19}}\\
        \vspace{10px}
        \textcolor{logo}{\hrulefill}\\
        \vspace{40px}
        \textbf{\Large Informazioni sul documento}\\
        \vspace{20px}
        \begin{tabular}{p{100px} | p{100px}}
            \textbf{Versione}     & 1.0.0                                                                      \\
            \textbf{Responsabile} & Tito Scutari                                                               \\
            \textbf{Redattore}    & Nicolò Giaccone                                                            \\
            \textbf{Verificatore} & Samuele De Simone                                                          \\
            \textbf{Destinatari}  & Prof. Tullio Vardanega \newline Prof. Riccardo Cardin, \newline TechSWEave \\
            \textbf{Stato}        & Approvato                                                                  \\
            \textbf{Uso}          & Interno                                                                    \\
        \end{tabular}\\
        \vspace{60px}
        \href{mailto:techsweave@gmail.com}{techsweave@gmail.com}\\

    \end{center}
\end{titlepage}

%INDICE
\newpage
\tableofcontents
\newpage

\section{Informazioni riunione}
\begin{itemize}
    \item Data e ora: 2021-04-19 ore 18:00-18:30.
    \item Luogo: Google Meet fornito dai proponenti.
    \item Partecipanti interni: Samuele De Simone, Manuel Varo, Tito Scutari, Nicolò Giaccone, Amedeo Meggiolaro, Marco Barbaresco.
    \item Partecipanti esterni: Milo Ertola, Alessandro Maccagnan.
    \item Segretario: Amedeo Meggiolaro.
\end{itemize}
\section{Argomenti della riunione}
\begin{itemize}
    \item Scelta identity manager: il team inizialmente aveva scelto di utilizzare Auth0 come identity manager, ma essendo questo deprecato per sistemi serverless, si è scelto, in accordo con il proponente di utilizzare Amazon Cognito;
    \item Scelta del framework: il team ha chiesto ai proponenti se si potesse utilizzare AWS Amplify dato che è un framework ben integrato con i servizi AWS, i proponenti lo hanno sconsigliato in quanto avrebbe portato ad onerosi costi di setup per l'ambiente di lavoro che avrebbero inciso sul tempo di codifica. Il gruppo ha quindi deciso di utilizzare il framework Serverless.
\end{itemize}
\section{Resoconto}
\begin{center}
    \textbf{\Large Resoconto riunione}\\
    \vspace{10px}
    \begin{table}[h!]
        \centering
        \rowcolors{2}{logo!10}{logo!40}
        \renewcommand{\arraystretch}{1.8}
        \begin{tabular}{p{150px} p{150px}}
            \rowcolor{logo!70} \textbf{Quesito} & \textbf{Soluzione}                                        \\
            Scelta identity manager             & Il gruppo ha scelto di utilizzare Amazon Cognito          \\
            Scelta del framework                & Il gruppo ha scelto di utilizzare il framework Serverless \\
        \end{tabular}
    \end{table}
\end{center}

\end{document}